\chapter{Nominal ECEF Mechanization}

\section{Corrected IMU Increment Inputs}

The high-rate mechanization consumes the bias-corrected increments
$\Delta\hat{\boldsymbol\theta}_{ib}^{b}$ and
$\Delta\hat{\mathbf{v}}_{ib}^{b}$ defined by the input contract in
\cref{sec:imu_input_contract}. These estimated bias states are part of the
nominal state and must be used in the strapdown integration loop, not only in
the lower-rate Kalman update path \cite{groves2013,titterton2004}.

\section{Coning and Sculling Compensation}

For a two-sample update over one mechanization interval, let
$(\Delta\hat{\boldsymbol\theta}_1,\Delta\hat{\mathbf{v}}_1)$ and
$(\Delta\hat{\boldsymbol\theta}_2,\Delta\hat{\mathbf{v}}_2)$ be the two
bias-corrected sub-sample increments, where sample 1 precedes sample 2 within
the interval $[t_k,t_k+\Delta t]$. The first v1 implementation should use the
standard two-sample coning correction
\begin{equation}
\Delta\boldsymbol\Theta
=
\Delta\hat{\boldsymbol\theta}_1
+
\Delta\hat{\boldsymbol\theta}_2
+
\frac{2}{3}
\Delta\hat{\boldsymbol\theta}_1
\times
\Delta\hat{\boldsymbol\theta}_2,
\label{eq:v1_coning_two_sample}
\end{equation}
and a two-sample sculling-corrected specific-force velocity increment
\begin{align}
\Delta\mathbf{V}^{b_{\mathrm{mid}}}
&=
\Delta\hat{\mathbf{v}}_1
+
\Delta\hat{\mathbf{v}}_2
+
\frac{2}{3}
\left(
\Delta\hat{\boldsymbol\theta}_1
\times
\Delta\hat{\mathbf{v}}_2
+
\Delta\hat{\mathbf{v}}_1
\times
\Delta\hat{\boldsymbol\theta}_2
\right).
\label{eq:v1_sculling_two_sample}
\end{align}
The increment in \cref{eq:v1_sculling_two_sample} is resolved in the
interval-midpoint body frame. The mean-rotation contribution is supplied by
the midpoint body-to-ECEF attitude in
\cref{eq:v1_discrete_velocity_update}; therefore the implementation must not
also add
$\tfrac{1}{2}(\Delta\hat{\boldsymbol\theta}_1+
\Delta\hat{\boldsymbol\theta}_2)\times
(\Delta\hat{\mathbf{v}}_1+\Delta\hat{\mathbf{v}}_2)$ to this increment.
Doing both would double count the first-order mean body rotation. For a
single raw increment, $\Delta\mathbf{V}^{b_{\mathrm{mid}}}$ is the measured
increment interpreted at the interval midpoint.

The exact coefficient convention must be verified against the selected IMU
increment ordering in tests. Higher-order multi-sample algorithms are deferred
until the two-sample path is correct and profiled. The v1 formulas are the
standard two-sample coning correction and the corresponding midpoint-frame
form of the two-sample sculling correction used for incremental IMU
mechanization \cite{groves2013,titterton2004}.

If the selected IMU reports internally compensated output intervals, the
Navigator does not combine two of those intervals with
\cref{eq:v1_coning_two_sample,eq:v1_sculling_two_sample}. It advances them
individually, retaining the same midpoint-body delta-velocity contract. The
compile-time simulator/product handshake selects one compensation owner.

\section{Quaternion Attitude Propagation}

For continuous-time notation, the relative angular rate of the body frame with
respect to ECEF, resolved in body, is
\begin{equation}
\omegab_{eb}^{b}
=
\hat{\omegab}_{ib}^{b}
-
C_e^b\omegab_{ie}^{e}.
\end{equation}
In the discrete implementation, $\Delta\boldsymbol\Theta$ from
\cref{eq:v1_coning_two_sample} represents the corresponding body-frame
increment over the interval after Earth-rate handling.

For v1, evaluate Earth rate with the attitude at the beginning of the
mechanization interval:
\begin{equation}
\Delta\boldsymbol\Theta_{eb,k}^{b}
=
\Delta\boldsymbol\Theta_{ib,k}^{b}
-
\hat{C}_{e,k}^{b}\omegab_{ie}^{e}\Delta t,
\qquad
\hat{C}_{e,k}^{b}=\left(\hat{C}_{b,k}^{e}\right)^T.
\label{eq:v1_discrete_earth_rate_attitude_increment}
\end{equation}
Here $\Delta\boldsymbol\Theta_{ib,k}^{b}$ is the coning-corrected body wrt
inertial increment from \cref{eq:v1_coning_two_sample}, and
$\Delta\boldsymbol\Theta_{eb,k}^{b}$ is the body wrt ECEF increment used for
the quaternion update.

Convert the corrected rotation vector to a scalar-first quaternion with
\cref{eq:v1_rotation_vector_quaternion}:
\begin{equation}
\delta\quat_{eb,k}^{b}
=
q\left(\Delta\boldsymbol\Theta_{eb,k}^{b}\right).
\label{eq:v1_discrete_attitude_delta_quat}
\end{equation}

Let a scalar-first quaternion use product $\otimes$ and let
$\boldsymbol\omega_q=[0,\ \omegab_{eb}^{bT}]^T$. The body-to-ECEF quaternion is
propagated with right multiplication by the relative angular-rate quaternion:
\begin{equation}
\boxed{
\dot{\quat}_{b}^{e}
=
\frac{1}{2}
\quat_b^e
\otimes
\boldsymbol\omega_q.
}
\label{eq:v1_quat_propagation}
\end{equation}
The corresponding v1 discrete quaternion update is
\begin{equation}
\boxed{
\hat{\quat}_{b,k+1}^{e,-}
=
\hat{\quat}_{b,k}^{e,+}
\otimes
\delta\quat_{eb,k}^{b}.
}
\label{eq:v1_discrete_quat_update}
\end{equation}
As a sign sanity check, if $\hat{\quat}_{b,k}^{e,+}=[1,0,0,0]^T$,
$\omegab_{ie}^{e}=\zero$, and
$\Delta\boldsymbol\Theta_{eb,k}^{b}=[0,0,\epsilon]^T$ for small positive
$\epsilon$, then
\begin{equation}
\hat{\quat}_{b,k+1}^{e,-}
\approx
\begin{bmatrix}
1 & 0 & 0 & \epsilon/2
\end{bmatrix}^{T},
\label{eq:v1_quaternion_sign_sanity}
\end{equation}
before normalization. This simple case pins the scalar-first quaternion
ordering, the right-multiplication propagation side, and the positive
body-frame yaw increment convention used by the implementation tests.
After discrete propagation, normalize:
\begin{equation}
\hat{\quat}_{b,k+1}^{e,-}
\leftarrow
\frac{\hat{\quat}_{b,k+1}^{e,-}}{\norm{\hat{\quat}_{b,k+1}^{e,-}}}.
\label{eq:v1_discrete_quat_normalization}
\end{equation}

\section{Velocity and Position Propagation}

The configured v1 gravity policy returns mass-attraction gravitation
$\boldsymbol{\gamma}^{e}$, not plumb-bob gravity. Define the outward
centrifugal pseudo-acceleration and plumb-bob gravity by
\begin{align}
\mathbf{a}_{\mathrm{cf}}^{e}(\mathbf{p}^{e})
&=
-\omegab_{ie}^{e}\times
\left(\omegab_{ie}^{e}\times\mathbf{p}^{e}\right),\\
\gb^{e}(\mathbf{p}^{e})
&=
\boldsymbol{\gamma}^{e}(\mathbf{p}^{e})
+
\mathbf{a}_{\mathrm{cf}}^{e}(\mathbf{p}^{e}).
\label{eq:v1_apparent_gravity_definition}
\end{align}

Using plumb-bob gravity $\gb^{e}(\pb_{eb}^{e})$, the ECEF velocity equation is
\begin{equation}
\boxed{
\dot{\vb}_{eb}^{e}
=
C_b^e\hat{\fb}_{ib}^{b}
+
\gb^{e}\left(\pb_{eb}^{e}\right)
-
2\omegab_{ie}^{e}\times\vb_{eb}^{e}.
}
\label{eq:v1_velocity_mechanization}
\end{equation}
The corresponding position equation is
\begin{equation}
\boxed{
\dot{\pb}_{eb}^{e}
=
\vb_{eb}^{e}.
}
\end{equation}

The mass-attraction vector and its gradient used by the mechanization and
covariance prediction come from the configured gravity policy. The rotating
planet policy supplies the centrifugal term and its exact gradient. For Earth
products, the intended v1 baseline is a WGS-84/J2-capable mass-attraction model
rather than a hard-coded spherical approximation \cite{groves2013}.

For implementation, the configured gravity policy must provide both
\begin{equation}
\boldsymbol{\gamma}_k^e
=
\boldsymbol{\gamma}^e\left(\hat{\pb}_{eb,k}^{e}\right),
\qquad
G_{\gamma,k}^{e}
=
\left.
\frac{\partial\boldsymbol{\gamma}^e}{\partial\pb_{eb}^{e}}
\right|_{\hat{\pb}_{eb,k}^{e}}.
\label{eq:v1_gravity_policy_contract}
\end{equation}
The mechanization constructs
\begin{equation}
\gb_k^e
=
\boldsymbol{\gamma}_k^e+\mathbf{a}_{\mathrm{cf},k}^e,
\qquad
G_{p,k}^{e}
=
G_{\gamma,k}^{e}
-
\Omega_{ie}^{e}\Omega_{ie}^{e},
\label{eq:v1_apparent_gravity_gradient}
\end{equation}
where
$-\Omega_{ie}^{e}\Omega_{ie}^{e}
=\partial\mathbf{a}_{\mathrm{cf}}^e/\partial\mathbf{p}^e$.
The covariance prediction uses this complete apparent-gravity gradient.

\section{Discrete Velocity and Position Update}

Using the midpoint-body-resolved sculling-corrected specific-force increment
$\Delta\mathbf{V}_{ib,k}^{b_{\mathrm{mid}}}$ from
\cref{eq:v1_sculling_two_sample}, define
\begin{align}
\delta\quat_{\mathrm{mid},k}
&=
q\left(\frac{1}{2}\Delta\boldsymbol\Theta_{eb,k}^{b}\right),\\
\hat{\quat}_{b,\mathrm{mid},k}^{e}
&=
\operatorname{normalize}
\left(
\hat{\quat}_{b,k}^{e,+}\otimes\delta\quat_{\mathrm{mid},k}
\right),\\
\bar C_{b,k}^{e}
&=
C_b^e\left(\hat{\quat}_{b,\mathrm{mid},k}^{e}\right).
\label{eq:v1_midpoint_attitude}
\end{align}
This normalized midpoint quaternion is the v1 approximation to the exact
interval-average body-to-ECEF transform. The v1 velocity update is
\begin{equation}
\boxed{
\hat{\vb}_{eb,k+1}^{e,-}
=
\hat{\vb}_{eb,k}^{e,+}
+
\bar C_{b,k}^{e}\Delta\mathbf{V}_{ib,k}^{b_{\mathrm{mid}}}
+
\left(
\gb_k^e
-
2\omegab_{ie}^{e}\times\hat{\vb}_{eb,k}^{e,+}
\right)\Delta t.
}
\label{eq:v1_discrete_velocity_update}
\end{equation}
The v1 position update uses trapezoidal velocity integration:
\begin{equation}
\boxed{
\hat{\pb}_{eb,k+1}^{e,-}
=
\hat{\pb}_{eb,k}^{e,+}
+
\frac{1}{2}
\left(
\hat{\vb}_{eb,k}^{e,+}
+
\hat{\vb}_{eb,k+1}^{e,-}
\right)
\Delta t.
}
\label{eq:v1_discrete_position_update}
\end{equation}
The superscript $+$ denotes the post-correction state entering the propagation
interval, and $-$ denotes the predicted state after propagation.

For covariance prediction over the same interval, use the interval-average
specific force corresponding to the same sculling-corrected velocity increment:
\begin{equation}
\boxed{
\hat{\fb}_{ib,k}^{b}
=
\frac{\Delta\mathbf{V}_{ib,k}^{b_{\mathrm{mid}}}}{\Delta t}.
}
\label{eq:v1_interval_average_specific_force}
\end{equation}
This midpoint-body-frame quantity is used to form
$\hat{\fb}_{ib,k}^{e}=\bar C_{b,k}^{e}\hat{\fb}_{ib,k}^{b}$ in the
interval-specific $F_k$ matrix. The same midpoint transform is used for
accelerometer/gyro bias coupling and continuous-noise mapping during this v1
interval linearization.

\section{Discrete Integration}

The first implementation should state:
\begin{itemize}
    \item the IMU increment ordering used by coning/sculling compensation;
    \item where biases are subtracted;
    \item whether Earth rate is evaluated at the beginning, midpoint, or end of
    the interval;
    \item when the quaternion is normalized;
    \item how gravity and the gravity gradient are evaluated.
\end{itemize}

\begin{implbox}
The concrete mechanization may expose an internal \texttt{integrate} operation.
The Navigator-level policy remains a broader propagation/time-advance seam
because the same step may also coordinate filter prediction, state-transition
history, and aiding-measurement buffering.
\end{implbox}
